% Distributed shared memory: processes on three nodes read and write one
% shared address space; every page of it resides in the memory that one of
% the nodes contributes, and remote accesses cross the network.
% Usage: % Distributed shared memory: processes on three nodes read and write one
% shared address space; every page of it resides in the memory that one of
% the nodes contributes, and remote accesses cross the network.
% Usage: % Distributed shared memory: processes on three nodes read and write one
% shared address space; every page of it resides in the memory that one of
% the nodes contributes, and remote accesses cross the network.
% Usage: % Distributed shared memory: processes on three nodes read and write one
% shared address space; every page of it resides in the memory that one of
% the nodes contributes, and remote accesses cross the network.
% Usage: \input{figures/l15-dsm}   (width ~118 mm)
\begin{tikzpicture}[x=1mm, y=1mm, font=\scriptsize\sffamily,
  >={Stealth[length=1.3mm]},
  mach/.style={draw, rounded corners=0.8mm, fill=white},
  proc/.style={draw, circle, fill=white, inner sep=0pt, minimum size=4.8mm},
  dsm/.style={draw, fill=ln-box, minimum width=30mm, minimum height=4.4mm, inner sep=0pt},
  frm/.style={draw, minimum width=8.4mm, minimum height=4.6mm, inner sep=0pt},
  cell/.style={draw, minimum width=13mm, minimum height=5mm, inner sep=0pt},
  lab/.style={font=\scriptsize\sffamily, text=ln-muted, inner sep=0.8pt},
  who/.style={font=\scriptsize\sffamily\bfseries, anchor=north west, inner sep=1pt}]
  \colorlet{lnNodeA}{ln-accent!22}
  \colorlet{lnNodeB}{ln-tip!22}
  \colorlet{lnNodeC}{ln-caution!16}
  % shared address space
  \node[lab, anchor=south west] at (1,40.8)
    {\iflangja{共有アドレス空間（すべてのプロセスから見える）}{shared address space (seen by every process)}};
  \foreach \ci/\cc in {0/lnNodeA,1/lnNodeB,2/lnNodeC,3/lnNodeA,4/lnNodeB,5/lnNodeC,6/lnNodeA,7/lnNodeB,8/lnNodeC} {
    \node[cell, fill=\cc] at ({7.5+13*\ci},38) {$p_{\ci}$};
  }
  % nodes
  \foreach \nk/\nx/\ncol/\npa/\npb/\fa/\fb/\fc in
      {1/0/lnNodeA/1/2/0/3/6, 2/41/lnNodeB/3/4/1/4/7, 3/82/lnNodeC/5/6/2/5/8} {
    \draw[mach] (\nx,0) rectangle ++(36,26);
    \node[who] at (\nx+0.5,25.5) {\iflangja{ノード}{node} \nk};
    \node[proc] at (\nx+15,19.5) {P\npa};
    \node[proc] at (\nx+27,19.5) {P\npb};
    \draw[<->] (\nx+15,21.9) -- (\nx+15,35.4);
    \draw[<->] (\nx+27,21.9) -- (\nx+27,35.4);
    \node[dsm] at (\nx+18,12) {DSM};
    \foreach \fj/\fp in {0/\fa,1/\fb,2/\fc} {
      \node[frm, fill=\ncol] at (\nx+8.5+9.5*\fj,4.6) {$p_{\fp}$};
    }
    \draw[thick, ln-muted] (\nx+18,0) -- (\nx+18,-5);
  }
  \node[lab, anchor=west] at (27.6,30.5) {\iflangja{読み書き}{read / write}};
  % network
  \draw[line width=1.2pt, ln-muted] (0,-5) -- (118,-5);
  \node[lab, anchor=north] at (59,-5.8)
    {\iflangja{ネットワーク（RDMA 対応のインターコネクトなど）}{network (e.g.\ an interconnect with RDMA)}};
\end{tikzpicture}
   (width ~118 mm)
\begin{tikzpicture}[x=1mm, y=1mm, font=\scriptsize\sffamily,
  >={Stealth[length=1.3mm]},
  mach/.style={draw, rounded corners=0.8mm, fill=white},
  proc/.style={draw, circle, fill=white, inner sep=0pt, minimum size=4.8mm},
  dsm/.style={draw, fill=ln-box, minimum width=30mm, minimum height=4.4mm, inner sep=0pt},
  frm/.style={draw, minimum width=8.4mm, minimum height=4.6mm, inner sep=0pt},
  cell/.style={draw, minimum width=13mm, minimum height=5mm, inner sep=0pt},
  lab/.style={font=\scriptsize\sffamily, text=ln-muted, inner sep=0.8pt},
  who/.style={font=\scriptsize\sffamily\bfseries, anchor=north west, inner sep=1pt}]
  \colorlet{lnNodeA}{ln-accent!22}
  \colorlet{lnNodeB}{ln-tip!22}
  \colorlet{lnNodeC}{ln-caution!16}
  % shared address space
  \node[lab, anchor=south west] at (1,40.8)
    {\iflangja{共有アドレス空間（すべてのプロセスから見える）}{shared address space (seen by every process)}};
  \foreach \ci/\cc in {0/lnNodeA,1/lnNodeB,2/lnNodeC,3/lnNodeA,4/lnNodeB,5/lnNodeC,6/lnNodeA,7/lnNodeB,8/lnNodeC} {
    \node[cell, fill=\cc] at ({7.5+13*\ci},38) {$p_{\ci}$};
  }
  % nodes
  \foreach \nk/\nx/\ncol/\npa/\npb/\fa/\fb/\fc in
      {1/0/lnNodeA/1/2/0/3/6, 2/41/lnNodeB/3/4/1/4/7, 3/82/lnNodeC/5/6/2/5/8} {
    \draw[mach] (\nx,0) rectangle ++(36,26);
    \node[who] at (\nx+0.5,25.5) {\iflangja{ノード}{node} \nk};
    \node[proc] at (\nx+15,19.5) {P\npa};
    \node[proc] at (\nx+27,19.5) {P\npb};
    \draw[<->] (\nx+15,21.9) -- (\nx+15,35.4);
    \draw[<->] (\nx+27,21.9) -- (\nx+27,35.4);
    \node[dsm] at (\nx+18,12) {DSM};
    \foreach \fj/\fp in {0/\fa,1/\fb,2/\fc} {
      \node[frm, fill=\ncol] at (\nx+8.5+9.5*\fj,4.6) {$p_{\fp}$};
    }
    \draw[thick, ln-muted] (\nx+18,0) -- (\nx+18,-5);
  }
  \node[lab, anchor=west] at (27.6,30.5) {\iflangja{読み書き}{read / write}};
  % network
  \draw[line width=1.2pt, ln-muted] (0,-5) -- (118,-5);
  \node[lab, anchor=north] at (59,-5.8)
    {\iflangja{ネットワーク（RDMA 対応のインターコネクトなど）}{network (e.g.\ an interconnect with RDMA)}};
\end{tikzpicture}
   (width ~118 mm)
\begin{tikzpicture}[x=1mm, y=1mm, font=\scriptsize\sffamily,
  >={Stealth[length=1.3mm]},
  mach/.style={draw, rounded corners=0.8mm, fill=white},
  proc/.style={draw, circle, fill=white, inner sep=0pt, minimum size=4.8mm},
  dsm/.style={draw, fill=ln-box, minimum width=30mm, minimum height=4.4mm, inner sep=0pt},
  frm/.style={draw, minimum width=8.4mm, minimum height=4.6mm, inner sep=0pt},
  cell/.style={draw, minimum width=13mm, minimum height=5mm, inner sep=0pt},
  lab/.style={font=\scriptsize\sffamily, text=ln-muted, inner sep=0.8pt},
  who/.style={font=\scriptsize\sffamily\bfseries, anchor=north west, inner sep=1pt}]
  \colorlet{lnNodeA}{ln-accent!22}
  \colorlet{lnNodeB}{ln-tip!22}
  \colorlet{lnNodeC}{ln-caution!16}
  % shared address space
  \node[lab, anchor=south west] at (1,40.8)
    {\iflangja{共有アドレス空間（すべてのプロセスから見える）}{shared address space (seen by every process)}};
  \foreach \ci/\cc in {0/lnNodeA,1/lnNodeB,2/lnNodeC,3/lnNodeA,4/lnNodeB,5/lnNodeC,6/lnNodeA,7/lnNodeB,8/lnNodeC} {
    \node[cell, fill=\cc] at ({7.5+13*\ci},38) {$p_{\ci}$};
  }
  % nodes
  \foreach \nk/\nx/\ncol/\npa/\npb/\fa/\fb/\fc in
      {1/0/lnNodeA/1/2/0/3/6, 2/41/lnNodeB/3/4/1/4/7, 3/82/lnNodeC/5/6/2/5/8} {
    \draw[mach] (\nx,0) rectangle ++(36,26);
    \node[who] at (\nx+0.5,25.5) {\iflangja{ノード}{node} \nk};
    \node[proc] at (\nx+15,19.5) {P\npa};
    \node[proc] at (\nx+27,19.5) {P\npb};
    \draw[<->] (\nx+15,21.9) -- (\nx+15,35.4);
    \draw[<->] (\nx+27,21.9) -- (\nx+27,35.4);
    \node[dsm] at (\nx+18,12) {DSM};
    \foreach \fj/\fp in {0/\fa,1/\fb,2/\fc} {
      \node[frm, fill=\ncol] at (\nx+8.5+9.5*\fj,4.6) {$p_{\fp}$};
    }
    \draw[thick, ln-muted] (\nx+18,0) -- (\nx+18,-5);
  }
  \node[lab, anchor=west] at (27.6,30.5) {\iflangja{読み書き}{read / write}};
  % network
  \draw[line width=1.2pt, ln-muted] (0,-5) -- (118,-5);
  \node[lab, anchor=north] at (59,-5.8)
    {\iflangja{ネットワーク（RDMA 対応のインターコネクトなど）}{network (e.g.\ an interconnect with RDMA)}};
\end{tikzpicture}
   (width ~118 mm)
\begin{tikzpicture}[x=1mm, y=1mm, font=\scriptsize\sffamily,
  >={Stealth[length=1.3mm]},
  mach/.style={draw, rounded corners=0.8mm, fill=white},
  proc/.style={draw, circle, fill=white, inner sep=0pt, minimum size=4.8mm},
  dsm/.style={draw, fill=ln-box, minimum width=30mm, minimum height=4.4mm, inner sep=0pt},
  frm/.style={draw, minimum width=8.4mm, minimum height=4.6mm, inner sep=0pt},
  cell/.style={draw, minimum width=13mm, minimum height=5mm, inner sep=0pt},
  lab/.style={font=\scriptsize\sffamily, text=ln-muted, inner sep=0.8pt},
  who/.style={font=\scriptsize\sffamily\bfseries, anchor=north west, inner sep=1pt}]
  \colorlet{lnNodeA}{ln-accent!22}
  \colorlet{lnNodeB}{ln-tip!22}
  \colorlet{lnNodeC}{ln-caution!16}
  % shared address space
  \node[lab, anchor=south west] at (1,40.8)
    {\iflangja{共有アドレス空間（すべてのプロセスから見える）}{shared address space (seen by every process)}};
  \foreach \ci/\cc in {0/lnNodeA,1/lnNodeB,2/lnNodeC,3/lnNodeA,4/lnNodeB,5/lnNodeC,6/lnNodeA,7/lnNodeB,8/lnNodeC} {
    \node[cell, fill=\cc] at ({7.5+13*\ci},38) {$p_{\ci}$};
  }
  % nodes
  \foreach \nk/\nx/\ncol/\npa/\npb/\fa/\fb/\fc in
      {1/0/lnNodeA/1/2/0/3/6, 2/41/lnNodeB/3/4/1/4/7, 3/82/lnNodeC/5/6/2/5/8} {
    \draw[mach] (\nx,0) rectangle ++(36,26);
    \node[who] at (\nx+0.5,25.5) {\iflangja{ノード}{node} \nk};
    \node[proc] at (\nx+15,19.5) {P\npa};
    \node[proc] at (\nx+27,19.5) {P\npb};
    \draw[<->] (\nx+15,21.9) -- (\nx+15,35.4);
    \draw[<->] (\nx+27,21.9) -- (\nx+27,35.4);
    \node[dsm] at (\nx+18,12) {DSM};
    \foreach \fj/\fp in {0/\fa,1/\fb,2/\fc} {
      \node[frm, fill=\ncol] at (\nx+8.5+9.5*\fj,4.6) {$p_{\fp}$};
    }
    \draw[thick, ln-muted] (\nx+18,0) -- (\nx+18,-5);
  }
  \node[lab, anchor=west] at (27.6,30.5) {\iflangja{読み書き}{read / write}};
  % network
  \draw[line width=1.2pt, ln-muted] (0,-5) -- (118,-5);
  \node[lab, anchor=north] at (59,-5.8)
    {\iflangja{ネットワーク（RDMA 対応のインターコネクトなど）}{network (e.g.\ an interconnect with RDMA)}};
\end{tikzpicture}
